% Source: https://artofproblemsolving.com/community/c3h2111244_equilateral_triangle_problem
% Stable identity: aops-topic-2111244-q01
Rotate the plane by $60^\circ$ about $B$ so that $A$ goes to $C$, and let $P'$ be the image of $P$.

\begin{asy}
usepackage("amsmath");
size(5.5cm);
defaultpen(fontsize(10pt));
real s=sqrt(25+12*sqrt(3));
pair B=(0,0), C=(s,0), A=(s/2,s*sqrt(3)/2);
real px=(16-25+s^2)/(2*s); pair P=(px,sqrt(16-px^2));
pair Q=rotate(-60,B)*P;
draw(A--B--C--cycle,linewidth(0.9));
draw(P--A); draw(P--B); draw(P--C);
draw(B--Q--C,blue+dashed+linewidth(0.8)); draw(P--Q,red+dashed+linewidth(0.8));
pair e1=unit(P-Q), e2=unit(C-Q);
draw(Q+0.35*e1--Q+0.35*e1+0.35*e2--Q+0.35*e2,red);
dot(A); dot(B); dot(C); dot(P); dot(Q,blue);
label("$A$",A,N); label("$B$",B,SW); label("$C$",C,SE); label("$P$",P,W); label("$P'$",Q,S,blue);
label("$3$",(P+A)/2,E); label("$4$",(P+B)/2,SE); label("$5$",(P+C)/2,NE);
label("$4$",(B+Q)/2,S,blue); label("$3$",(Q+C)/2,S,blue); label("$4$",(P+Q)/2,E,red);
label("\mbox{}",(min(currentpicture,user=true)+max(currentpicture,user=true))/2,black); // keeps the text after the figure black when LaTeX typesets the labels inline
\end{asy}

The rotation preserves lengths, so $BP'=BP=4$ and $P'C=PA=3$. Since $BP=BP'$ and $\angle PBP'=60^\circ$, triangle $PBP'$ is equilateral, so $PP'=4$ and $\angle BP'P=60^\circ$. In triangle $PP'C$ the sides are $PP'=4$, $P'C=3$, $PC=5$, and $3^2+4^2=5^2$, so by the converse of the Pythagorean theorem $\angle PP'C=90^\circ$.

The angle $\angle BP'C$ is therefore either $60^\circ+90^\circ=150^\circ$ (if $P$ lies inside angle $BP'C$, as in the figure) or $90^\circ-60^\circ=30^\circ$. By the law of cosines in triangle $BP'C$, with side $s=BC$:
\[s^2=4^2+3^2-2\cdot4\cdot3\cos\angle BP'C=25\mp12\sqrt3 .\]
The value $25-12\sqrt3\approx4.2$ is impossible, because a point inside an equilateral triangle is closer than $s$ to every vertex, so $s>PC=5$ and $s^2>25$. Hence $s^2=25+12\sqrt3$ and
\[s=\boxed{\sqrt{25+12\sqrt3}}\approx6.77 .\]
